\chapter{Conclusion\label{cha:conclusion}}

Sequential optimization of operator placement and push-pull communication in hierarchical fog-cloud Complex Event Processing systems fails to exploit interdependencies between placement and communication decisions, yielding suboptimal evaluation plans. This thesis presented \emph{Kraken}, a \emph{joint optimization framework} that integrates INEv placement costs with PrePP communication costs through a unified state-space search model addressing cost function integration (RQ1) and state-space formulation (RQ2). The framework employs \emph{dual plan evaluation} where both all-push and optimized push-pull communication strategies are computed and compared for \emph{every candidate placement} before selecting the globally optimal configuration. We implemented Kraken within the INES simulation environment and evaluated it across 21,717 experimental runs against four baselines—Central all-push, INEv placement-only, PrePP communication-only, and INES sequential—varying network topologies, query workloads, and event distribution parameters.

Kraken transmits \emph{only 1.6\% of baseline events} (median transmission ratio 0.016), undercutting the state-of-the-art INES sequential baseline by \emph{2.1$\times$} (INES: 0.033), PrePP by 3.2$\times$ (PrePP: 0.049), and INEv by \emph{32$\times$} (INEv: 0.500) (RQ3, Section~\ref{sec:rq3_comparison}). These cost reductions accompany 12\% superior latency performance compared to INES (median ratio 2.333 vs. 2.667) and 4.7$\times$ faster computation times (Kraken median: 0.277 seconds vs. INES median: 1.295 seconds). Across all 21,717 experimental runs, the framework demonstrates practical tractability with median computation time of 1.03 seconds where \emph{98.87\% of runs complete within one minute} (RQ4, Section~\ref{sec:rq4_tractability}), validating that heuristic search effectively navigates the exponentially large $O(m^k)$ solution space (Section~\ref{sec:correctness}) despite each placement evaluation requiring NP-complete PrePP optimization~\cite{purtzel_predicate-based_2022}. We evaluated multiple search strategies where Beam-3 captures the exploration-efficiency sweet spot, reducing transmission costs 7\% relative to Greedy search at 2.6$\times$ computational overhead (RQ5, Section~\ref{sec:search_strategy_impact}). Multi-objective weight configurations achieve flexible cost-latency trade-offs: latency-prioritized deployments ($\alpha = 0.3$) reduce latency by 50\% while still undercutting INES baseline costs, whereas balanced configurations ($\alpha = 0.6$) achieve 16\% latency reduction with only 4\% cost overhead (RQ6, Section~\ref{sec:multi_objective_tradeoffs}). Even with Local Event Correction (Section~\ref{sec:local_event_correction}) and Sink Cost Adjustment (Section~\ref{sec:sink_cost_adjustment}) disabled, Kraken outperforms INES (median 0.032 vs. 0.033), confirming that \emph{joint search inherently discovers superior solutions unavailable to sequential approaches}.

The framework operates under abstract cost models optimizing transmission volume (events per window) and network distance (hops) rather than physical resources including bandwidth capacity, packet serialization overhead, CPU utilization, or memory constraints. Three critical limitations define future research directions: INEv's all-push split function prunes the search space before joint optimization, potentially eliminating placement-communication combinations optimal only under push-pull costs; the framework evaluates only the single minimal-cost PrePP plan per placement candidate, ignoring alternative plans with different cost-latency profiles; and static assumptions exclude dynamic adaptation to node failures, link congestion, and time-varying event rates (Section~\ref{sec:correctness}). Despite these simplifications, Kraken demonstrates consistent superiority across diverse experimental conditions, validating that joint optimization methodology generalizes robustly beyond specific deployment scenarios.

Future work should optimize the split function through memoization and aggressive pruning, extend the framework to evaluate multiple PrePP plans per placement candidate to fully explore cost-latency trade-offs, and investigate \emph{multi-node placement strategies} beyond the current single-node placement constraint. Practical deployment requires integrating physical resource models with dynamic adaptation mechanisms, while advanced search strategies including A* with admissible heuristics could further improve solution quality. Kraken establishes that \emph{joint optimization consistently and substantially outperforms sequential approaches} for hierarchical fog-cloud CEP systems, providing a validated foundation for distributed stream processing optimization where placement and communication decisions exhibit strong interdependencies.